\documentclass[a4paper,11pt]{article}

\usepackage{titlesec}
\usepackage{color}
\usepackage{enumitem}
\usepackage{fancyhdr}
\usepackage{latexsym}
\usepackage[empty]{fullpage}
\usepackage[hidelinks]{hyperref}
\usepackage[normalem]{ulem}
\usepackage[english]{babel}
\usepackage{fontawesome5}

\input glyphtounicode
\pdfgentounicode=1

\usepackage[default]{sourcesanspro}
\urlstyle{same}

% Optimized margins for one page
\pagestyle{fancy}
\fancyhf{}
\renewcommand{\headrulewidth}{0in}
\renewcommand{\footrulewidth}{0in}
\setlength{\tabcolsep}{0in}
\addtolength{\oddsidemargin}{-0.6in}
\addtolength{\topmargin}{-0.8in}
\addtolength{\textwidth}{1.2in}
\addtolength{\textheight}{1.6in}
\raggedbottom{}


% Compact section styles
\titleformat{\section}
  {\scshape\normalsize\bfseries}{}
  {0em}{\color{blue}}[\color{black}\titlerule\vspace{-3pt}]
\titlespacing{\section}{0pt}{4pt}{3pt}

% Compact custom commands
\newcommand{\resumeItem}[1]{\item\small{#1}\vspace{-2pt}}
\newcommand{\resumeItemListStart}{\begin{itemize}[leftmargin=0.12in, label={\tiny$\bullet$}, itemsep=1pt, parsep=0pt, topsep=0pt]}
\newcommand{\resumeItemListEnd}{\end{itemize}\vspace{-3pt}}

\begin{document}

% Compact header
\begin{center}
\textbf{\large MBASSI EWOLO Loïc Aron}\\
\small{\textit{Recherche d'un stage de 4 mois et 3 semaines à partir d'avril 2026 en Génie Logiciel}}\\
\small{7 rue les chênes d'or, 95000, Cergy-Pontoise, France}\\
\small{\href{loicaron.mbassiewolo@ensea.fr}{ \faEnvelope \uline{loicaron.mbassiewolo@ensea.fr}} | \faPhone (+33) 7 59 52 33 98 | \href{https://www.linkedin.com/in/mbassi-loic-3942a9246/}{\faLinkedin \ @mbassi-loic} |
\href{https://github.com/Nameless0l}{\faGithub \ @Nameless0l} |
\href{https://mbassiloic.tech/#portfolio}{\faBriefcase  \ portfolio} |
\href{https://medium.com/@wwwmbassiloic}{\faMedium \ medium}
\href{https://mbassiloic.tech/#blog}{\faBlog \ blog}
\end{center}
\vspace{-11pt}


\section{Compétences}
\resumeItemListStart
\item \textbf{Technique - IA/ML :} TensorFlow, PyTorch, Scikit-learn, Keras, OpenCV, MediaPipe | \textbf{Déploiement sur Docker}
\item \textbf{Programmation :} Python, Java, JavaScript, TypeScript, HTML/CSS | \textbf{Outils :} Git, VS Code, MySQL, PostgreSQL, Figma
\item \textbf{Frameworks :} React.js, NextJS, Spring Boot \textbf{| Science des Données :} Pandas, NumPy, Matplotlib, Jupyter
\item \textbf{Compétences Transversales :} Travail d'équipe, Créativité, Pensée critique
\item \textbf{Langues :} Français (Langue maternelle), Anglais (Communication scientifique)
\resumeItemListEnd

\section{Formation}
\resumeItemListStart
\item \textbf{École nationale supérieure de l'électronique et de ses applications (ENSEA)} \hfill 2025-Présent\\
2ème année - Formation en SIA (Signal et IA), Ingénieur de conception
\item \textbf{École Nationale Supérieure Polytechnique, Yaoundé} \hfill 2021-2025\\
-Niveau 3 \& 4 (Génie Informatique), compétences : Python pour ML, Développement full-stack, Analyse de données\\
-Niveau 1 \& 2 - compétences : Logique, algorithmique, programmation en C | Licence 1 en Informatique en parallèle
\resumeItemListEnd

\section{Expérience Professionnelle}
\resumeItemListStart
\item \textbf{Stage à distance chez MILA, Québec, Canada} \hfill 06/2025 - 08/2025
\begin{itemize}[leftmargin=0.12in, itemsep=-1pt]
\item Grokking (d'OpenAI), Généralisation des réseaux de neurones après sur-apprentissage, test des travaux sur des MLP 
\end{itemize}
\item \textbf{Freelance, à distance - CSC, Douala} \hfill 10/2024 - 08/2025
\begin{itemize}[leftmargin=0.12in, itemsep=-1pt]
\item FinTech, développement d'une API et des interfaces de gestion financières pour des petites entreprises 
\end{itemize}
\item \textbf{Stagiaire - Laboratoire IOTA, ENSPY} \hfill 06/2024 - 09/2024
\begin{itemize}[leftmargin=0.12in, itemsep=-1pt]
\item Implémentation d'algorithmes pour la traduction de la langue des signes camerounaise
\end{itemize}
\resumeItemListEnd


\section{Projets Significatifs}
\resumeItemListStart

\item \textbf{Upgrade Jetson - Challenge CoHoMa}  \hfill Présent - académique
\begin{itemize}[leftmargin=0.12in, itemsep=-3pt]
\vspace{-2pt}
\item Mise en service des Jetsons pour des robots de l'armée, intégration et optimisation des modèles d'IA légers
\end{itemize}
\item \textbf{Harmony Gloves - Traduction de la langue des signes} \hfill 2024-Présent
\begin{itemize}[leftmargin=0.12in, itemsep=-3pt]
\vspace{-2pt}
\item Système basé sur des capteurs avec les réseaux de neurones pour la reconnaissance en temps réel des gestes
\end{itemize}
\vspace{-2pt}

\item \textbf{Yemba Language Learning Tool avec modèles d'IA pré-entraînés - Audio Langue locale au Cameroun } \hfill 2024
\begin{itemize}[leftmargin=0.12in, itemsep=-2pt]
\item Application développée en Python pour l'apprentissage des langues locales : \href{https://yemba-learning-audio.mbassiloic.tech/}{ici}
\end{itemize}

\vspace{-2pt}
\item \textbf{Système Multi-Agents basé sur l'ADK de Google - Optimisation de l'agriculture} \hfill 2024, projet personnel
\begin{itemize}[leftmargin=0.12in, itemsep=-2pt]
\item Système conçu pour optimiser les rendements agricoles au Cameroun grâce à 04 agents d'IA collaborant ensemble 
\end{itemize}

\item \textbf{Smart bin : solution intégrée éco-responsable avec IA et analyse statistique} \hfill Octobre - Novembre 2024
\begin{itemize}[leftmargin=0.12in, itemsep=-2pt]
\item Système IoT de poubelle intelligente basé sur les algorithmes de ML (Random Forest) et du voyageur de commerce pour l'optimisation de la collecte des déchets, \textbf{1er Prix} National CITS
\end{itemize}
\resumeItemListEnd
\vspace{-2pt}

\section{Activités Académiques}
\resumeItemListStart
\item \textbf{Cellule IA ENSPY} - Chef de la cellule \hfill 2024-2025 | Préparation aux hackathons, veille technologique, leader des projets
\item \textbf{Chef cellule Projet Informatique} \hfill 2023, 2024 | Réalisation des projets informatiques de l'école | Management
\item \textbf{Vice-Président club Génie Informatique} -  \hfill 2023, 2024 | Réalisation et encadrements de certains projets tech de l'école (Projet antivirus) et management des équipes et cellules du club 
\item \textbf{MLPC (Machine Learning Project Competition)} \hfill 2023, 2024, 2025 | IA pour l'analyse des ECG, détection d'expressions misogynes
\resumeItemListEnd

\section{Réalisations et Récompenses}
\resumeItemListStart
\item \textbf{Hackathons et Récompenses Majeures :} 1er prix Hackathon CITS (Oct 2024) | 2ème prix national hackathon d'Indabax (Challenge IA d'Afrique) | 1er prix Hackathon Douala (Août 2024) | 2ème prix Semaine de la Fille Scientifique (Mars 2023) | 1er prix journée de la fille scientifique (Mars 2023) | Nuit de l'informatique France 2024 (Déc 2024)
\item \textbf{Publications :} Mini \href{https://mbassiloic.tech/#blog}{\uline{blogs}} sur l'IA | Présentations d'articles et débats sur les thèmes d'IA (Ex : risque existentiel)
\resumeItemListEnd

\section{Certifications}
\resumeItemListStart
\item \textbf{En cours :} Supervised ML: Regression \& Classification (DeepLearning.ai) | TOEIC | Agents IA en embarqué (IBM)
\item \textbf{Complétées :} ML Supervisé : Régression \& Classification (DeepLearning.ai) | Python for Data Science | POO Java (EPFL)
\resumeItemListEnd

\section{Références}
\resumeItemListStart
\item \textbf{Ing. James} - Ingénieur de Recherche Google DeepMind, Chef de la Cellule IA ENSPY | \href{mailto:moudie@deepmind.com}{\uline{moudie@deepmind.com}}
\item \textbf{Ing, PhD. Pascal} - Mila, Ingénieur de Recherche | \href{mailto:pascalnotsawo@gmail.com}{\uline{pascalnotsawo@gmail.com}}
\item \textbf{Dr. Anne Marie CHANA} - Directrice du Laboratoire IOTA | \href{mailto:anne.chana@univ-yaounde1.cm}{\uline{anne.chana@univ-yaounde1.cm}}
\resumeItemListEnd

\end{document}